\section{实际验证、证据范围与复现}

前面几节给出了从图到递归有限密度、解析证书和采样分布的连续构造。
本节说明哪些步骤已经实际执行，以及每种核查支持什么结论。
项目采用完整的有质量零动量 Taylor 条件，BPH 的 Bogoliubov 递归始终是主线；
尚未核验的一般 $Z$ 部分不作为数值结果或证明的前提。

\subsection{698 张图的覆盖意味着什么}

一至四圈、外腿数一至三的固定标签图共 698 张。全部图都已产生可执行表达式，
并在相等正参数点求值；每张图仍保存原始边索引，平行边可以分别进入子图。
全部规范图与精确对称权重也已经逐项匹配官方 feyngen，而不只匹配总数。
一、两圈保留旧图号；高圈图号采用新的精确染色规范序。
其他生成器中的图号只有在核对带固定外腿的同构关系后才能对应。

这里的完整覆盖是\emph{图枚举和表达式编译覆盖}。
698 个内部参数点都能求值，不等于 698 个完整振幅都已有高精度积分，
更不等于所有 UV 面附近的误差和原动量收敛都已统一证明。
这些层次在目录、执行记录和解析证书中分别保存。

\subsection{44 项独立有限密度检查}

独立参考程序从约化关联矩阵求路由，不复用生产程序的树叶消元；
它以 mpmath 构造 Gaussian 矩阵，再对边界缩放变量直接高精度微分，
不调用生产 Taylor jet、矩阵消元或密度算术。
它保留混合二次系数，包括多个嵌套二点反项的乘积。
另用精确生成树与生成二森林多项式检查原始 Gaussian 项的 $U,F$，
范围包含选定三圈、四圈 primitive 图。

实际复跑的 44 项检查全部通过，覆盖一、两圈完整图集的多个参数点，
以及高圈的选定 primitive、嵌套二点反项和 UV 条件控制实例。
这里的真嵌套由各检查图自身边序下的严格包含直接确认：三圈检查图的
$\{0,1\}\subsetneq\{0,1,2,3,8\}$；四圈检查图的
$\{0,1\}$ 与 $\{3,4\}$ 分别严格包含于 $\{0,1,2,3,4,5,6,7\}$。
它们分别对应新目录的 \texttt{L3E3G53} 与 \texttt{L4E3G300}，
各检查四种参数点模式；这与下述完整积分选用的图分开记录。
其中固定级数截断反例另用 160、200 位直接微分复核，
采用独立的 $10^{-80}$ 绝对阈值，避免一般相对容差掩盖大密度中的偏差。
对该点，生产 Decimal200 与高精度参考的差约为
$3.53\times10^{-122}$。
这些检查证实列出的代数与数值条件；它们不是整个 simplex 的积分误差包络。

\subsection{14 项完整两圈积分互校}

联合 UV 算术修复后，通用编译器已在当前冻结源码下重新实际执行
14 个两圈图与运动学输入的完整积分，
再与之前独立推导并核验的有限 BPH 公式比较。
专用公式只在验证工具中作为参考调用；通用生产计算仍由输入图自动生成。
全部 14 项比较通过，实际最大绝对差为
\[
 \max |J_{\rm generic}-J_{\rm reference}|
 =6.009245532560825\times10^{-6}.
\]
两圈验证工具先从完整UV子图判断分层提案是否适用。适用图使用固定总尺度形状$1/2$、
内部Beta形状3和bare边比例形状1，从12组各$2^{14}$点开始，
只有自身标准误不大于$5\times10^{-6}$且半样本差不大于$10^{-5}$才停止加密；
其余图保持原19次幂、12组的通用提案。参考值在选阶后才调用。
实际逐次加密历史与采样总量保存在结果中。此前低效普通提案的中断运行单独保留，
没有计作完成14项，也没有从屏幕打印伪造原始采样记录。
参考方案、每圈测度、质量平方、固定外腿和有限 Taylor 条件保持一致。
这一步比较的是完整有限积分，而不只是极点、$U,F$ 或某个内部参数点。
该最大差描述这 14 次运行；不能把它推广为所有运动学输入的严格统一精度。
独立随机组标准误、半样本细化差与实际参考差均保留在结果 JSON 中。

\subsection{35 张高圈 primitive 图的解析证书}

三圈 4 张、四圈 31 张 primitive 三点图，共 35 张，
都已实际构造图导出的 Hepp 子集证书。
证书检查全部真边子集的正 UV gap，并以精确有理动态规划计算 Hepp 系数。
上一节证明把该系数连接到参数绝对可积上界、原始动量 $L^1$ 控制，
以及共同 $D_\eta$ 正则化的移除；质量和外动量的适用范围随证书明确记录。

这 35 张是\emph{高圈 primitive 三点图}的解析覆盖，不能解释为所有高圈图已闭合。
对于有真发散子图的三圈、四圈图，裸图周期界不能直接控制完整递归余项。
采样时将非正 gap 截成正值只是定义重要性分布，不会自动成为递归 UV 证明。
一般的递归 UV 余项估计及其对内部子图外动量的统一控制仍有待完成。

\subsection{一次实际原生三圈有限 BPH 互校}

官方 pySecDec 1.6.6 已为 \texttt{L3E3G5} 实际生成并编译 C++ 原生库。
该图是新三圈三点目录的第 5 张图，包含 7 个顶点、9 条内部边，
物理外腿标签 $1,2,3$ 分别接在顶点 $4,5,6$，没有真发散子图，
其 Hepp 系数为 192。上游生成 300 个 sector；生成耗时约 32.99 秒，
原生编译耗时约 193.37 秒。随后分别在共同 $M=1$ 的零点和单位对称点，
以不同种子实际运行原生 QMC，耗时约 107.23 秒和 103.11 秒。

这里 $I=9,L=3,d=6-2\epsilon$。含完整 Gamma 前因子的参数表达式为
\[
 I_{\rm DR}(p,M)=\Gamma(3\epsilon)
       \int_\Delta U^{-3+4\epsilon}F(p,M)^{-3\epsilon}.
\]
令 $P=\int_\Delta U^{-3}$，其共同极点是 $P/(3\epsilon)$。
若 $c_0$ 表示完整有限系数，则
\[
 c_0(p,M)=\int_\Delta U^{-3}
       \left[\frac43\log U-\log F(p,M)-\gamma_E\right].
\]
由于图中没有真发散子图，BPH 递归只剩整体零阶减除。
在 $M=1$、$F(0,1)=U$ 下，严格有限转换正是
\[
 J_R(p,1)=c_0(p,1)-c_0(0,1)
 =-\int_\Delta U^{-3}\log\frac{F(p,1)}U.
\]
没有额外的 3 或 $1/3$ 因子；也没有乘入图权重、耦合或 MS-bar 因子。
因此这里比较的是同一个有限 BPH 条件，不能把裸 MS 极点直接充当有限振幅。

实际结果为
\begin{align*}
 c_0(0,1)&=-0.9987759652197846\ \pm3.25\times10^{-6},\\
 c_0((1,1,1),1)&=-1.0721646501342081\ \pm4.45\times10^{-6},\\
 J_R^{\rm native}((1,1,1),1)
 &=-0.07338868491442352\ \pm5.51\times10^{-6}.
\end{align*}
两次极点估计差为 $1.03\times10^{-6}$，其组合报告标准差为
$1.26\times10^{-6}$；极点重建的周期也与独立正生成树多项式采样一致。
通用图编译器使用 Hepp 采样、16 组各 $2^{22}$ 个点，得到
\[
 J_R^{\rm generic}=-0.07338791561208803\ \pm6.19\times10^{-6},
\]
半样本细化差为 $6.64\times10^{-6}$，两条方法的实际绝对差为
$7.69\times10^{-7}$。上述正负误差都是随机数值诊断，不是严格区间。

原始级数、种子、输入图、脚本、生成与编译日志、库哈希及时间均已保留。
执行期间的网页渲染修复改变过全包指纹；完整旧源码快照证实数学依赖逐文件未变，
比较记录同时保留旧执行指纹和当前指纹，没有覆盖原始变化标志。
联合 UV 算术修复后，上述通用 Hepp QMC 已在当前冻结源码下重新实际执行；
上游原始级数重解析不被称为又一次原生积分，旧通用 QMC 仍保留历史执行记录。

\subsection{真嵌套三圈图的完整积分试算}

另一次实际运行选择 \texttt{L3E3G53}，在目录边序中具有
$\{5,6\}\subsetneq\{0,1,3,5,6\}$ 两个严格嵌套的二点子图，
两者的减除阶数均为 $\omega=2$。
在 $M=1$、$(s_1,s_2,s_3)=(1,1,1)$ 下，
完整递归编译密度使用 8 组共 524288 个采样点，得到
\[
 J_R=-0.00023303604812\ \pm5.59\times10^{-5}.
\]
更换独立随机化及重要性参数，另用 8 组共 262144 个点，得到
$-0.00029981283515\ \pm2.52\times10^{-5}$。
两者差为 $6.68\times10^{-5}$，约为组合报告标准误的 1.09 倍；
两次半样本细化差分别为 $3.67\times10^{-5}$ 和 $1.91\times10^{-5}$。
这一结果证明程序已实际完成真嵌套图的完整积分试算，当前精度仍较粗。
两次计算复用同一图导出有限密度，属于独立随机化的一致性检查，
不是独立解析核互校，也没有给出严格积分误差界。

\subsection{四圈完整递归和的独立核对}

第\ref{sec:grouped-sampling}节的分层工具根据输入图识别三个不交bubble及收缩后的一圈结构，
完整密度仍由\texttt{compile\_bph}构造，不调用第\ref{sec:four-loop-independent}节的参考公式。
默认总尺度形状参数为1；经保留的pilot比较后，主核查明确使用总尺度形状$1/2$、
bubble内部Beta形状3及bare边比例形状1。
实际12组共196608点给出
\[
 J_{\rm generic}=-0.0005154702382511922\ \pm1.80923\times10^{-6}.
\]
独立外圈动量参考为$-0.0005158785412715456$，实际绝对差为$4.08303\times10^{-7}$；
主采样半样本细化差为$1.18853\times10^{-6}$。
独立种子8组共65536点给出
$-0.0005116648680529154\ \pm3.69832\times10^{-6}$，
两次随机化相差$3.80537\times10^{-6}$。
正负数值均指报告的随机组标准误，不是严格总误差区间。

普通Dirichlet及初始平面分组未达到精度的原始记录仍保留，
不会因最终提案达到目标而改写为成功。新的采样权重同时经过确定性单项式矩与独立PDF检查；
参考标量只在完成积分后比较。图号一致性修复后，主结果及独立种子已用正式工具重新实际执行，
工具起止SHA、完整输入图、程序SHA和冻结核心指纹均保存。
完整比较见\path{results/automatic_four_loop_independent.json}，
原始采样及失败对照见\path{docs/automatic-group-sampling-evidence/raw-files.json}。
该例提供一个实际四圈非primitive完整有限积分核对，并不代表所有四圈图均已精密求积。

\subsection{当前完成范围与未闭合部分}

现在已经有完整图编译、低圈完整有限积分互校、35 张高圈 primitive 解析证书，
以及一张三圈 primitive 的真正外部原生后端互校。
选定高圈实例的实际执行与完成状态单独记录。
其中 \texttt{L3E3G21} 与 \texttt{L4E3G205} 分别含两张、三张不交二点子图，
这两张图的完整积分检验的是多重不交二点反项，本身不构成真嵌套的证据；
真嵌套图的粗精度试算见上一小节，它仍缺少独立解析核互校。
四圈多重不交二点反项的完整独立核对见上一小节；选定高圈粗精度执行另见
\path{results/automatic_higher_integrals.json}。
点态 UV 检查不能替代整个参数积分控制，采样误差也不能替代原动量的绝对收敛证明。
698 张图的编译覆盖不等于 698 张图均已完成高精度全积分，
一般递归 UV 证明和任意圈归纳尚未闭合。

\subsection{复现入口与证据文件}

在解压后的软件目录运行。表达式构造和 primitive 证书只需 Python3.11 或更新版本；
通用积分使用可选 NumPy、SciPy，独立高精度微分参考需要 mpmath，
原生 Laurent 级数解析另需 pySecDec。
\begin{samepage}
\begin{verbatim}
python3 -m phi3lab build-automatic
python3 -m phi3lab automatic --loops 3 --legs 3 --index 5
python3 -m phi3lab automatic --input my-graph.json --integrate
python3 tools/compare_automatic_bph.py
python3 tools/verify_automatic_integrals.py
python3 tools/compare_automatic_native.py
python3 tools/run_automatic_higher.py
python3 tools/verify_automatic_four_loop.py
python3 tools/verify_tests.py
\end{verbatim}
\end{samepage}
自定义图 JSON 使用整数顶点字段 \texttt{vertices}、边列表 \texttt{edges}
和固定标签外腿列表 \texttt{external}，与目录中的图格式一致。
边的重排或顶点的重命名不会切换到按名字指定的专用积分公式。

\begin{itemize}
\item 图集与精确权重：\path{results/higher_graph_feyngen_comparison.json}。
\item 698 张图的表达式：\path{results/automatic/catalog.json} 及其逐图文件。
\item 44 项独立密度检查：\path{results/automatic_bph_comparison.json}。
\item 14 项完整两圈比较：\path{results/automatic_integral_comparison.json}。
\item 原生三圈有限转换与比较：\path{results/automatic_native_comparison.json}。
\item 真嵌套三圈完整积分试算：\path{results/automatic_nested_integral.json}。
\item 实际原生构建记录：\path{docs/automatic-native-build.json}。
\end{itemize}
每份证据的原始执行时间、源码指纹与实际工作阶段分别保存。
重新解析原始输出和重新执行积分是两个不同操作，不能混写。

\subsection{主要来源与方案边界}

历史 Henge 递归见 Kennedy、Binoth 与 Rippon，
\href{https://arxiv.org/abs/0712.1016}{arXiv:0712.1016}。
Brown 与 Kreimer 的
\href{https://arxiv.org/abs/1112.1180}{arXiv:1112.1180}
提供更广泛的参数重整化背景，但其 single-scale 构造没有被直接替换成本项目的完整有质量零动量条件。
现代 pySecDec 更新见 Heinrich 等，
\href{https://arxiv.org/abs/2305.19768}{arXiv:2305.19768}，
CPC 295 (2024) 108956；实际版本、输入和已执行点在各自记录中明确给出。
官方组合图参考为 \url{https://github.com/michibo/feyncop}。
Tropical 积分参考
\href{https://arxiv.org/abs/2302.08955}{arXiv:2302.08955}
要求合适的 quasi-finite 输入，不能据此直接把有子发散裸图当成已匹配 BPH 的参考。
本项目的递归编译器、Taylor jets 与采样代码由协作实现，没有打包第三方科学源码；
本文不宣称方法新颖性或一般圈数定理已经完成。
